% !TEX spellcheck = en_US
% !TEX spellcheck = LaTeX
\documentclass[letterpaper,10pt,english]{article}
\usepackage{%
amsfonts,%
amsmath,%
amssymb,%
amsthm,%
babel,%
bbm,%
%biblatex,%
caption,%
centernot,%
color,%
enumerate,%
%enumitem,%
epsfig,%
epstopdf,%
etex,%
fancybox,%
framed,%
fullpage,%
%geometry,%
graphicx,%
hyperref,%
latexsym,%
mathptmx,%
mathtools,%
multicol,%
paralist,%
pgf,%
pgfplots,%
pgfplotstable,%
pgfpages,%
proof,%
psfrag,%
%subfigure,%
tcolorbox,%
tfrupee,%
tikz,%
times,%
%ulem,%
url,%
xcolor,%
mathpazo,%
}

\definecolor{shadecolor}{gray}{.95}%{rgb}{1,0,0}
\usepackage[margin=1in,top=0.75in]{geometry}
\usepackage[mathscr]{eucal}
\usepgflibrary{shapes}
\usepgfplotslibrary{fillbetween}
\usetikzlibrary{%
arrows,%
backgrounds,%
chains,%
decorations.pathmorphing,% /pgf/decoration/random steps | erste Graphik
decorations.text,% 
matrix,%
positioning,% wg. " of "
fit,%
patterns,%
petri,%
plotmarks,%
scopes,%
shadows,%
shapes.misc,% wg. rounded rectangle
shapes.arrows,%
shapes.callouts,%
shapes%
}

%\pgfplotsset{compat=newest} %<------ Here
\pgfplotsset{compat=1.11} %<------ Or use this one

\theoremstyle{plain}
\newtheorem{thm}{Theorem}[section]
\newtheorem{lem}[thm]{Lemma}
\newtheorem{prop}[thm]{Proposition}
\newtheorem{cor}[thm]{Corollary}
\newtheorem{clm}[thm]{Claim}

\theoremstyle{definition}
\newtheorem{defn}[thm]{Definition}
\newtheorem{conj}[thm]{Conjecture}
\newtheorem{exmp}[thm]{Example}
\newtheorem{axiom}[thm]{Axiom}
\newtheorem{assum}[thm]{Assumption}
\newtheorem{exerc}[thm]{Exercise}
\newtheorem*{defin}{Definition}

\theoremstyle{remark}
\newtheorem{rem}{Remark}
\newtheorem{note}{Note}

\newcommand{\iid}{\emph{i.i.d.}\ }
\newcommand{\Cov}{\operatorname{Cov}}
%\newcommand{\det}{\operatorname{det}}
%\newcommand{\Real}{\mathbb{R}}
\newcommand{\tr}{\operatorname{tr}}
\newcommand{\Var}{\operatorname{Var}}
\newcommand{\cov}{\operatorname{cov}}

%\renewcommand{\proof}[1]{\begin{proof}#1\end{proof}}
\newcommand{\EQ}[1]{\begin{equation*}#1\end{equation*}}
\newcommand{\EQN}[1]{\begin{equation}#1\end{equation}}
\newcommand{\eq}[1]{\begin{align*}#1\end{align*}}
\newcommand{\meq}[2]{\begin{xalignat*}{#1}#2\end{xalignat*}}
\newcommand{\norm}[1]{\left\lVert#1\right\rVert}
\newcommand{\inner}[1]{\left\langle #1\right\rangle}
\newcommand{\abs}[1]{\left\lvert#1\right\rvert}
\newcommand{\expect}[1]{\mathbb{E}\left[{#1}\right]}
\newcommand{\prob}[1]{\mathbb{P}\left[{#1}\right]}
\newcommand{\given}{\; \big\vert \;} 
\newcommand{\set}[1]{\left\{#1\right\}} 
\newcommand{\indicator}[1]{1_{\set{#1}}} 
\newcommand{\Ind}[1]{\mathbbm{1}_{#1}} 
\newcommand{\SetIn}[1]{\mathbbm{1}_{\set{#1}}} 
\newcommand{\red}[1]{\textcolor{red}{#1}} 
\newcommand{\floor}[1]{\left\lfloor#1\right\rfloor}
\newcommand{\ceil}[1]{\left\lceil#1\right\rceil}


\newcommand{\C}{\mathbb{C}}
\newcommand{\D}{\mathbb{D}}
\newcommand{\E}{\mathbb{E}}
\newcommand{\F}{\mathbb{F}}
\newcommand{\N}{\mathbb{N}}
\renewcommand{\P}{\mathbb{P}}
\newcommand{\Q}{\mathbb{Q}}
\newcommand{\R}{\mathbb{R}}
\newcommand{\Z}{\mathbb{Z}}

\newcommand{\bA}{\mathbf{A}}
\newcommand{\bI}{\mathbf{I}}
\newcommand{\bU}{\mathbf{1}}
\newcommand{\bx}{\mathbf{x}}
\newcommand{\bX}{\mathbf{X}}
\newcommand{\bZ}{\mathbf{Z}}


\newcommand{\cA}{\mathcal{A}}
\newcommand{\cB}{\mathcal{B}}
\newcommand{\cC}{\mathcal{C}}
\newcommand{\cF}{\mathcal{F}}
\newcommand{\cG}{\mathcal{G}}
\newcommand{\cM}{\mathcal{M}}
\newcommand{\cN}{\mathcal{N}}
\newcommand{\cP}{\mathcal{P}}
\newcommand{\cT}{\mathcal{T}}
\newcommand{\cX}{\mathcal{X}}
\newcommand{\cY}{\mathcal{Y}}

\newcommand{\sA}{\mathscr{A}}
\newcommand{\sB}{\mathscr{B}}
\newcommand{\sC}{\mathscr{C}}
\newcommand{\sE}{\mathscr{E}}
\newcommand{\sF}{\mathscr{F}}
\newcommand{\sG}{\mathscr{G}}
\newcommand{\sH}{\mathscr{H}}
\newcommand{\sL}{\mathscr{L}}
\newcommand{\sS}{\mathscr{S}}
\newcommand{\sT}{\mathscr{T}}
\newcommand{\sX}{\mathscr{X}}
\newcommand{\sY}{\mathscr{Y}}
\newcommand{\sZ}{\mathscr{Z}}

\newcommand{\tS}{\tilde{S}}
% Debug
\newcommand{\todo}[1]{\begin{color}{blue}{{\bf~[TODO:~#1]}}\end{color}}

% a few handy macros

\renewcommand{\le}{\leqslant}
\renewcommand{\ge}{\geqslant}
\newcommand\matlab{{\sc matlab}}
\newcommand{\goto}{\rightarrow}
\newcommand{\bigo}{{\mathcal O}}
%\newcommand{\half}{\frac{1}{2}}
%\newcommand\implies{\quad\Longrightarrow\quad}
\newcommand\reals{{{\rm l} \kern -.15em {\rm R} }}
\newcommand\complex{{\raisebox{.043ex}{\rule{0.07em}{1.56ex}} \hskip -.35em {\rm C}}}


% macros for matrices/vectors:

% matrix environment for vectors or matrices where elements are centered
\newenvironment{mat}{\left[\begin{array}{ccccccccccccccc}}{\end{array}\right]}
\newcommand\bcm{\begin{mat}}
\newcommand\ecm{\end{mat}}

% matrix environment for vectors or matrices where elements are right justifvied
\newenvironment{rmat}{\left[\begin{array}{rrrrrrrrrrrrr}}{\end{array}\right]}
\newcommand\brm{\begin{rmat}}
\newcommand\erm{\end{rmat}}

% for left brace and a set of choices
%\newenvironment{choices}{\left\{ \begin{array}{ll}}{\end{array}\right.}
\newcommand\when{&\text{if~}}
\newcommand\otherwise{&\text{otherwise}}
% sample usage:
%\delta_{ij} = \begin{choices} 1 \when i=j, \\ 0 \otherwise \end{choices}


% for labeling and referencing equations:
\newcommand{\eql}{\begin{equation}\label}
\newcommand{\eqn}[1]{(\ref{#1})}
% can then do
%\eql{eqnlabel}
%...
%\end{equation}
% and refer to it as equation \eqn{eqnlabel}.
% some useful macros for finite difference methods:
\newcommand\unp{U^{n+1}}
\newcommand\unm{U^{n-1}}
% for chemical reactions:
\newcommand{\react}[1]{\stackrel{K_{#1}}{\rightarrow}}
\newcommand{\reactb}[2]{\stackrel{K_{#1}}{~\stackrel{\rightleftharpoons}
 {\scriptstyle K_{#2}}}~}
\makeatletter
\def\th@plain{%
\thm@notefont{}% same as heading font
\itshape % body font
}
\def\th@definition{%
\thm@notefont{}% same as heading font
\normalfont % body font
}
\makeatother
\date{}


\title{Lecture-11: Generating functions}
\author{}

\begin{document}
\maketitle

%%%%%%%%%%%%%%%%%%%%%%%%%%%%%%%%%%%%%%%%%%%%%%%%%%%%%
\section{Generating functions}
%%%%%%%%%%%%%%%%%%%%%%%%%%%%%%%%%%%%%%%%%%%%%%%%%%%%%
Suppose that $X:\Omega \to \R$ is a continuous random variable on the probability space $(\Omega,\sF, P)$ with distribution function $F_X: \R \to [0,1]$. 

%\begin{lem}
%If moment generating function $\Phi_X$ exists for a random variable $X:\Omega \to \R$ all $u \in \R$, 
%then all moments exist, and $\E X^k = \frac{d^k\Phi_X(u)}{du^k}\rvert_{u=0}$. 
%\end{lem}
%
%\begin{defn}[Characteristic function] 
%For a random variable $X$, the \textbf{characteristic function} $\Phi_X$ is defined by $\Phi_X(u) \triangleq \E e^{j u X}$, where $j = \sqrt{-1}$. 
%\end{defn}

%%%%%%%%%%%%%%%%%%%%%%%%%%%%%%%%%%%%%%%%%%%%%%%%%%%%%%%%
\subsection{Characteristic function}
%%%%%%%%%%%%%%%%%%%%%%%%%%%%%%%%%%%%%%%%%%%%%%%%%%%%%%%%
\begin{shaded}
\begin{exmp} 
%Consider a probability space $(\Omega, \sF, P)$ 
%and a random variable $X: \Omega \to \R$ defined on this probability space. 
Let $j \triangleq \sqrt{-1}$, 
then we can show that $h_u: \R \to \C$ defined by $h_u(x) \triangleq e^{ju x} = \cos(u x) + j \sin(u x)$ is also Borel measurable for all $u \in \R$.  
Thus, $h_u(X):\Omega \to \C$ is a complex valued random variable on this probability space.  
%Since $\abs{\cos(u X)} \le1$ and $\abs{\sin(u X)} \le 1$, the mean $\E h_u(X) = \E\cos(u X) + j \E\sin(u X)$ always exists and is finite. 
\end{exmp} 
\end{shaded}
%\begin{defn}[Characteristic function]
%For a random variable $X: \Omega \to \R$ defined on the probability space $(\Omega, \sF, P)$, 
%the \textbf{characteristic function} $\Phi_X: \R \to \C$ is defined by $\Phi_X(u) = \E e^{ju X}$ for all $u \in \R$.
%%, where $j = \sqrt{-1}$ and $e^{ju x} = \cos(u x) + \sin(u x)$. 
%\end{defn}
\begin{defn}%[Characteristic function] 
For a random variable $X: \Omega \to \R$ defined on the probability space $(\Omega, \sF, P)$, 
the \textbf{characteristic function} $\Phi_X: \R \to \C$ is defined by $\Phi_X(u) \triangleq \E e^{j u X}$ for all $u \in \R$ and $j^2 = -1$. 
%, where $j = \sqrt{-1}$ and $e^{ju x} = \cos(u x) + \sin(u x)$. 
\end{defn}
%\begin{lem}
%If moment generating function $\Phi_X$ exists for a random variable $X:\Omega \to \R$ all $u \in \R$, 
%then all moments exist, and $\E X^k = \frac{d^k\Phi_X(u)}{du^k}\rvert_{u=0}$. 
%\end{lem}
%
%\begin{defn}[Characteristic function] 
%For a random variable $X$, the \textbf{characteristic function} $\Phi_X$ is defined by $\Phi_X(u) \triangleq \E e^{j u X}$, where $j = \sqrt{-1}$. 
%\end{defn}
%\begin{defn} 
%Let $j \in \C$ such that $j^2 = -1$, then the \textbf{characteristic function} $\Phi_X: \R \to \R$, of a random variable $X$ is defined as 
%\EQ{
%\Phi_X(u) \triangleq \E[e^{juX}], \text{ for all }u \in \R.
%}
%\end{defn}
%\begin{rem}
%Since $e^{ju} = \cos{u} + j\sin{u}$ for real $u \in \R$, the characteristic function can be equivalently written as $\Phi_X(u) = \E[\cos{u X}] + j\E[\sin{u X}]$.
%\end{rem} 
\begin{rem}
The characteristic function $\Phi_X(u)$ is always finite, 
since $\abs{\Phi_X(u)} = \abs{\E e^{ju X}} \le \E\abs{e^{ju X}} = 1$. 
%\EQ{
 %\abs{\Phi_X{(u)}} = \abs{\int_{-\infty}^{\infty}e^{juX}f_X(x)dx}
 %\le \int_{-\infty}^{\infty}\abs{e^{juX}f_X(x)}dx= \int_{-\infty}^{\infty}f_X(x)dx =1.
 %}
\end{rem} 
\begin{rem}
For a discrete random variable $X: \Omega \to \sX$ with PMF $P_X:\sX\to [0,1]$, 
the characteristic function $\Phi_X(u) %= \E[e^{juX}] 
= \sum_{x \in \sX}e^{ju x}P_X(x).$
\end{rem} 
\begin{rem}
For a continuous random variable $X: \Omega \to \R$ with density function $f_X: \R \to \R_+$, 
the characteristic function $\Phi_X(u) = \int_{-\infty}^{\infty}e^{ju X}f_X(x)dx.$
\end{rem} 
%\begin{rem}
%Following statements are true for characteristic functions. 
%\begin{enumerate} 
%\item 
%\item 
%\item 
%\item 
%\end{enumerate}
%\end{rem}
%\begin{lem}
%If $\E[X^k]$ exists and is finite for an integer $k \in \N$, then the derivatives of $\Phi_X$ up to order $k$ exist and are continuous, and $\Phi_X^{(k)}(0) = j^k\E[X^k]$.
%\end{lem}
\begin{shaded}
\begin{exmp}[Gaussian random variable] 
For a Gaussian random variable $X: \Omega \to \R$ with mean $\mu$ and variance $\sigma^2$, 
the characteristic function $\Phi_X$ is 
\EQ{
\Phi_X(u) 
= \frac{1}{\sqrt{2\pi\sigma^2}}\int_{x\in\R}e^{ju x}e^{-\frac{(x-\mu)^2}{2\sigma^2}}dx
%= \frac{1}{\sqrt{2\pi\sigma^2}}\int_{x\in\R}\exp\Big(-\frac{(x^2+\mu^2-2\mu x-2ju\sigma^2x)}{2\sigma^2}\Big)\\
%= \frac{e^{\frac{(\mu+ju\sigma^2)^2-\mu^2}{2\sigma^2}}}{\sqrt{2\pi\sigma^2}}\int_{x\in\R}\exp\Big(-\frac{(x- \mu - ju\sigma^2)^2}{2\sigma^2}\Big).
= \exp\Big(-\frac{u^2\sigma^2}{2} + ju\mu\Big).
}
We observe that $\abs{\Phi_X(u)} = e^{-u^2\sigma^2/2}$ has Gaussian decay with zero mean and variance $1/\sigma^2$. 
\end{exmp}
\end{shaded}
%\begin{lem}
%If $\E[X^k]$ exists and is finite for an integer $k \in \N$, then the derivatives of $\Phi_X$ up to order $k$ exist and are continuous, and $\Phi_X^{(k)}(0) = j^k\E[X^k]$.
%\end{lem}
\begin{thm}
If $\E\abs{X}^N$ is finite for some integer $N \in \N$, 
then $\Phi_X^{(k)}(u)$ is finite and continuous functions of $u \in \R$ for all $k \in [N]$. 
Further, $\Phi_X^{(k)}(0) = j^k\E[X^k]$ for all $k \in [N]$.  
\end{thm}
\begin{proof}
%Let us differentiate the characteristic function $\Phi_X{(u)}$, with respect to $u$, to write
%\EQ{
%\Phi_X^{'}(u) = \frac{d\Phi_X(u)}{du}= \frac{d}{du}\left(\int_{-\infty}^{\infty}e^{juX}f_X(x)dx\right). 
%}
Since $\E\abs{X}^N$ is finite, then so is $\E\abs{X}^k$ for all $k\in [N]$. 
Therefore, $\E[X^k]$ exists and is finite. 
Exchanging derivative and the integration (which can be done since $e^{ju x}$ is a bounded function with all derivatives), 
and evaluating the derivative at $u=0$, we get 
$
\Phi_X^{(k)}(0) 
= \E\left[\frac{d^ke^{ju X}}{du^k}\Big\rvert_{u=0}\right]
%= \int_{-\infty}^{\infty}\frac{d^ke^{juX}}{du^k}\Big\rvert_{u=0}f_X(x)dx 
%= \int_{-\infty}^{\infty}j^kx^kf_X(x)dx 
%= \int_{-\infty}^{\infty}(j^kxe^{juX})\Big\rvert_{u=0}f_X(x)dx 
= j^k\E[X^k].
$
%Since $\E\abs{X}^N$ is finite, then so is $\E\abs{X}^k$ for all $k\in [N]$. 
%Therefore, $\E[X^k]$ exists and is finite, 
%and 
%%It turns out that $\Phi_X^{'}(0) = j\E[X]$, when both L.H.S and R.H.S are finite. 
%$\Phi_X^{(k)}(0) = j^k\E[X^k]$. 
\end{proof}

\begin{thm} 
Two random variables have the same distribution iff they have the same characteristic function. 
\end{thm}
\begin{proof}
It is easy to see the necessity and the sufficiency is difficult. 
\end{proof}

%%%%%%%%%%%%%%%%%%%%%%%%%%%%%%%%%%%%%%%%%%%%%%%%%%%%%%%%
\subsection{Moment generating function}
%%%%%%%%%%%%%%%%%%%%%%%%%%%%%%%%%%%%%%%%%%%%%%%%%%%%%%%%
\begin{shaded}
\begin{exmp} 
%Consider a probability space $(\Omega, \sF, P)$ 
%and a random variable $X: \Omega \to \R$ defined on this probability space. 
A function $g_t: \R \to \R_+$ defined by $g_t(x) \triangleq e^{t x}$ is monotone and hence Borel measurable for all $t \in \R$. 
Therefore, $g_t(X): \Omega \to \R_+$ is a positive random variable on this probability space.  
\end{exmp} 
\end{shaded}
%\begin{defn}[Moment generating function] 
%For a random variable $X: \Omega \to \R$ defined on the probability space $(\Omega, \sF, P)$, 
%the moment generating function $M_X: \R \to \R$ is defined by $M_X(u) = \E e^{u X}$ for all $u \in \R$ where $M_X(u)$ is finite. 
%\end{defn}
\begin{defn}%[Moment generating function] 
For a random variable $X: \Omega \to \R$ defined on the probability space $(\Omega, \sF, P)$, 
the \textbf{moment generating function} $M_X: \R \to \R_+$ is defined by $M_X(t) \triangleq \E e^{t X}$ for all $t \in \R$ where $M_X(t)$ is finite. 
\end{defn}
\begin{rem}
Characteristic function always exist, however are complex in general.  
Sometimes it is easier to work with moment generating functions, when they exist. 
\end{rem}
%
%\begin{defn} 
%For a random variable $X$, the \textbf{moment generating function} denoted by $M_X: \R \to \R$, is defined by
%\EQ{
%M_X(t)\triangleq \E[e^{tX}], \text{ for all } t \in \R.
%}
%\end{defn}
\begin{lem} 
For a random variable $X$, 
if the MGF $M_X(t)$ is finite for some $t \in \R$, 
then $M_X(t) = 1 + \sum_{n \in \N} \frac{t^n}{n!}\E[X^n].$
\end{lem}
\proof{ 
From the Taylor series expansion of $e^\theta$ around $\theta = 0$, we get $e^\theta = 1 +  \sum_{n \in \N} \frac{\theta^n}{n!}$. 
Therefore, taking $\theta = tX$, we can write 
$
e^{tX}= 1 +  \sum_{n \in \N} \frac{t^n}{n!}X^n.
$
Taking expectation on both sides, the result follows from the linearity of expectation, when both sides have finite expectation. 
}

\begin{shaded}
\begin{exmp}[Gaussian random variable] 
For a Gaussian random variable $X: \Omega \to \R$ with mean $\mu$ and variance $\sigma^2$, 
the moment generating function $M_X$ is 
$
M_X(t) = \exp\Big(\frac{t^2\sigma^2}{2}+t\mu\Big).
$
\end{exmp}
\end{shaded}

%%%%%%%%%%%%%%%%%%%%%%%%%%%%%%%%%%%%%%%%%%%%%%%%%%%%%%%%%%%
\subsection{Probability generating function}
%%%%%%%%%%%%%%%%%%%%%%%%%%%%%%%%%%%%%%%%%%%%%%%%%%%%%%%%%%%
For a non-negative integer-valued random variable $X: \Omega \to \sX \subseteq \Z_+$, 
it is often more convenient to work with the $z$-transform of the probability mass function, 
called the probability generating function. 
\begin{defn}
For a discrete non-negative integer-valued random variable $X: \Omega \to \sX \subseteq \Z_+$ with probability mass function $P_X: \sX \to [0,1]$, the \textbf{probability generating function} $\Psi_X: \C \to \C$ is defined by
\EQ{
\Psi_X(z) \triangleq \E[z^X] = \sum_{x \in \sX}z^xP_X(x),\quad z \in \C, \abs{z} \le 1. 
}
\end{defn} 
\begin{lem}
For a non-negative simple random variable $X: \Omega \to \sX$, we have $\abs{\Psi_X(z)} \le 1$ for all $\abs{z} \le 1$. 
%The absolute value of the probability generating function $\Psi_X$ evaluated at $z \in \C$ with $\abs{z} \le 1$ for a positive simple random variable $X$, is upper bounded by unity.
\end{lem}
\begin{proof}
Let $z \in \C$ with $\abs{z}\le 1$. 
Let $P_X : \sX \to [0,1]$ be the probability mass function of the non-negative simple random variable $X$. 
Since any realization $x \in \sX$ of random variable $X$ is non-negative, we can write 
\EQ{
\abs{\Psi_X(z)} = \abs{\sum_{x \in \sX}z^xP_X(x)} 
\le \sum_{x \in \sX}\abs{z}^x P_X(x) \le \sum_{x \in \sX} P_X(x)= 1.
}
\end{proof}
%\begin{defn}
%For a non-negative integer-valued random variable $X$ it is often more convenient
%to work with the $z$-transform of the PMF, defined by $\Psi_X(z) = \E z^X = \sum_{k \ge 0}z^kp_X(k)$, 
%for real or complex $z$ with $\abs{z} \le 1$.  
%\end{defn}

%\begin{thm} 
%Two  non-negative integer-valued random variables have the same probability distribution iff their $z$-transforms are equal. 
%If $\E[X^k]$ is finite it can be found from the derivatives of $\Psi_X$ up to the $k$th order at $z = 1$,
%$\Psi_X^{(k)}(1) = \E[X(X-1)\dots(X-k+1)]$. 
%\end{thm}
%\begin{proof}
%The necessity is clear. 
%For sufficiency, we see that $\Psi_X^{(k)}(0) = k!p_X(k)$. 
%Further, interchanging the derivative and the summation (by dominated convergence theorem), we get the second result. 
%\end{proof}


%\begin{defn} 
%For a positive random variable $X$, 
%the $k$-th order factorial moment of random variable $X$ is defined as 
%\EQ{
%\E\left[\prod_{i=0}^{k-1}(X-i)\right]=\E[X(X-1)(X-2)\dots(X-k+1)].
%}
%\end{defn}
\begin{thm} 
For a non-negative simple random variable $X:\Omega \to \sX$ with finite $k$th moment $\E X^k$, 
the $k$-th derivative of probability generating function evaluated at $z=1$ is the $k$-th order factorial moment of $X$. 
That is, 
\EQ{
\Psi_X^{(k)}(1) 
= \E\left[\prod_{i=0}^{k-1}(X-i)\right]
= \E[X(X-1)(X-2)\ldots(X-k+1)].
}
\end{thm}
\proof{
It follows from the interchange of derivative and expectation. 
%$\Psi_X^{'}(1) = \E[X]$ if both L.H.S and R.H.S are finite. In general, the $k$-th derivative is
}
\begin{rem}
Moments can be recovered from $k$th order factorial moments. 
For example, 
\meq{2}{
&\E[X] = \Psi_X^{'}(1), &&\E[X^2] = \Psi_X^{(2)}(1) + \Psi_X^{'}(1).
}
\end{rem}
\begin{thm} 
Two  non-negative integer-valued random variables have the same probability distribution iff their $z$-transforms are equal. 
%If $\E[X^k]$ is finite it can be found from the derivatives of $\Psi_X$ up to the $k$th order at $z = 1$,
%$\Psi_X^{(k)}(1) = \E[X(X-1)\dots(X-k+1)]$. 
\end{thm}
\begin{proof}
The necessity is clear. 
For sufficiency, we see that $\Psi_X^{(k)} = \sum_{x \ge k}k!z^{x-k}P_X(x)$ and hence $\Psi_X^{(k)}(0) = k!P_X(k)$. 
Further, interchanging the derivative and the summation (by dominated convergence theorem), we get the second result. 
\end{proof}
%%%%%%%%%%%%%%%%%%%%%%%%%%%%%%%%%%%%%%%%%%%%%%%%%%%%%%%%%%%%%%%%%%%
%\section{Gaussian Random Vectors}
%%%%%%%%%%%%%%%%%%%%%%%%%%%%%%%%%%%%%%%%%%%%%%%%%%%%%%%%%%%%%%%%%%%
%
%\begin{defn}
%For a random vector $X: \Omega\to \R^n$ defined on a probability space $(\Omega,\sF,P)$, 
%we can define the characteristic function $\Phi_X:\R^n\to\C$ by $\Phi_X(u) \triangleq \E e^{j \inner{u,X}}$ where $u \in \R^n$. 
%\end{defn}
%\begin{rem}
%If $X:\Omega\to\R^n$ is an independent random vector, 
%then $\Phi_X(u) = \prod_{i=1}^n\Phi_{X_i}(u_i)$ for all $u \in \R^n$.
%\end{rem}
%\begin{defn}%[\emph{I.i.d.} Gaussian random vector] 
%For a probability space $(\Omega, \sF, P)$, 
%an \textbf{\iid Gaussian random vector} $X: \Omega \to \R^n$ is a continuous random vector defined by its density function
%\EQ{
%f_X(x) = \frac{1}{(2\pi\sigma^2)^{n/2}}\exp\left(-\frac{1}{2}\sum_{i=1}^n\frac{(x_i-\mu_1)^2}{\sigma^2}\right)\text{ for all }x \in \R^n.
%}
%for some real scalar $\mu_1$ and positive $\sigma^2\in \R_+$.
%\end{defn}
%\begin{rem}
%For an \iid Gaussian random vector with density parametrized by $(\mu_1,\sigma^2)$, 
%the components are \iid Gaussian random variables with mean $\mu_1$ and variance $\sigma^2$.
%\end{rem}
%\begin{rem}
%The characteristic function $\Phi_X$ of an \iid Gaussian random vector $X:\Omega\to \R^n$ parametrized by $(\mu_1,\sigma^2)$ is given by 
%\EQ{
%\Phi_X(u) = \prod_{i=1}^n\Phi_{X_i}(u_i) = \exp\Big(-\frac{\sigma^2}{2}\sum_{i=1}^nu_i^2 +j\mu_1\sum_{i=1}^nu_i\Big).
%}
%\end{rem}
%\begin{lem} 
%For an \iid zero mean unit variance Gaussian vector $Z:\Omega\to\R^n$, vector $\alpha \in \R^n$, and scalar $\mu \in \R$, 
%the affine combination $Y \triangleq \mu + \inner{\alpha, Z}$ is a Gaussian random variable.  
%\end{lem}
%\begin{proof}
%From the linearity of expectation and the fact that $Z$ is a zero mean vector, 
%we get $\E Y = \mu$. 
%Further, from the linearity of expectation and the fact that $\E[ZZ^T] = I$, we get 
%\EQ{
%\sigma^2 \triangleq \Var(Y) = \E(Y-\mu)^2 = \sum_{i=1}^n\sum_{k=1}^n\alpha_i\alpha_k\E[Z_iZ_k] = \inner{\alpha,\alpha} = \norm{\alpha}^2_2 = \sum_{i=1}^n\alpha_i^2.
%}
%To show that $Y$ is Gaussian, 
%it suffices to show that $\Phi_Y(u) = \exp(-\frac{u^2\sigma^2}{2}+ju\mu)$ for any $u \in \R$. 
%Recall that $Z$ is an independent random vector with individual components being identically zero mean unit variance Gaussian. 
%Therefore, $\Phi_{Z_i}(u) = \exp(-\frac{u^2}{2})$, 
%and we can compute the characteristic function of $Y$ as 
%\EQ{
%\Phi_Y(u) 
%= \E e^{juY} 
%= e^{ju\mu}\E \prod_{i=1}^ne^{ju\alpha_iZ_i}
%= e^{ju\mu}\prod_{i=1}^n\Phi_{Z_i}(u\alpha_i)
%= \exp(-\frac{u^2\sigma^2}{2}+ju\mu).
%}
%\end{proof}
%\begin{defn}%[Gaussian random vector] 
%For a probability space $(\Omega, \sF, P)$, 
%a \textbf{Gaussian random vector} $X: \Omega \to \R^n$ can be written as 
%\EQ{
%X = \mu + AZ,
%}
%from some vector $\mu \in \R^n$, matrix $A \in \R^{n\times n}$, 
%and \iid Gaussian random vector $Z: \Omega \to \R^n$ with mean~$0$ and variance~$1$. 
%We denote the covariance matrix for the Gaussian vector $X$ by $\Sigma \triangleq \E(X-\mu)(X-\mu)^T$. 
%\end{defn}
%\begin{rem}
%The components of the Gaussian random vector are Gaussian random variables with mean $\mu_i$ and variance $\sum_{k=1}^nA_{i,k}^2 = (AA^T)_{i,i}$, 
%since each component $X_i = \mu_i + \sum_{k=1}^nA_{i,k}Z_k$ is an affine combination of zero mean unit variance \iid random variables. 
%\end{rem}
%\begin{lem}
%For a Gaussian random vector $X=\mu+AZ$ for $\mu\in\R^n, A \in \R^{n\times n}$, and \iid zero mean unit variance Gaussian random vector $Z$, 
%the covariance matrix is $\Sigma = AA^T$. 
%\end{lem}
%\begin{proof}
%We can write $X_i = \mu_i + \sum_{k=1}^nA_{i,k}Z_k$ and we get $\E X_i = \mu_i$ from linearity of expectations and the fact that $\E Z_k =0$ for all $k \in [n]$. 
%Similarly, the $(i,j)$th component of covariance matrix is the mean of 
%\EQ{
%(X_i-\mu_i)(X_j-\mu_j) 
%= \sum_{\ell=1}^n\sum_{k=1}^nA_{i,k}A_{j,\ell}Z_kZ_\ell
%= \sum_{k=1}^nA_{i,k}A_{j,k}Z_k^2 + \sum_{k\neq\ell}^nA_{i,k}A_{j,\ell}Z_kZ_\ell.
%}
%From the linearity of expectation, 
%and the fact that $Z$ is an independent zero mean unit variance random vector, 
%we get $\Sigma_{i,j} = (AA^T)_{i,j}$. 
%\end{proof}
%\begin{prop} 
%The density for a Gaussian random vector $X:\Omega\to\R^n$ with mean $\mu \in \R^n$ and invertible covariance matrix $\Sigma\in\R^{n \times n}$, 
%is given by 
%\EQ{
%f_X(x) = \frac{1}{\sqrt{(2\pi)^n\det(\Sigma)}}\exp\left(-\frac{1}{2}(x-\mu)^T\Sigma^{-1}(x-\mu)\right)\text{ for all }x \in \R^n.
%}
%%where $\mu \in \R^n$ and $\Sigma \in \R^{n \times n}$ is a positive definite matrix. 
%\end{prop}
%\begin{proof}
%We can write $X = \mu + \Sigma^{\frac{1}{2}}Z$, where $Z:\Omega\to\R^n$ is an \iid zero mean unit variance Gaussian random vector. 
%Then, we observe that $Z = \Sigma^{-\frac{1}{2}}(X-\mu)$. 
%This implies that the Jacobian matrix $J(z) = \Sigma^{-\frac{1}{2}}$, 
%since the $(i,j)$th component of the Jacobian matrix $J(z)$ is given by 
%$J_{ij}(z) = \frac{\partial z_j}{\partial x_i} = \Sigma^{-\frac{1}{2}}_{j,i},\quad i,j\in [n].$
%Recall that the density of $Z$ is $f_Z(z) = \frac{1}{\sqrt{(2\pi)^n}}\exp(-\frac{1}{2}z^Tz)$, 
%and from the transformation of random vectors, we get 
%\EQ{
%f_X(x) 
%= f_Z(\Sigma^{-\frac{1}{2}}(x-\mu))\det(\Sigma^{-\frac{1}{2}})
%= \frac{1}{(2\pi)^{n/2}\det(\Sigma)^{1/2}}\exp\Big(-\frac{1}{2}(x-\mu)^T\Sigma^{-1}(x-\mu)\Big),\quad x\in \R^n.
%}  
%\end{proof}
%\begin{rem}
%%To show the given form of the density $f_X$, we utilize the almots sure uniqueness of inverse of characteristic functions. 
%%We show that a direct computation of characteristic function $\Phi_X$ matches to the one obtained by the given form of density function. 
%For any $u \in \R^n$, 
%we compute the characteristic function $\Phi_X$ from the distribution of $Z$ as 
%\EQ{
%\Phi_X(u) 
%= \E e^{j\inner{u,X}} 
%= \E\exp\Big(j\inner{u,\mu} + j\inner{A^Tu,Z}\Big)
%= \exp(j\inner{u,\mu})\Phi_Z(A^Tu) 
%= \exp(j\inner{u,\mu}-\frac{1}{2}u^T\Sigma u).
%}
%%Assuming the density of $X$ as given in the Theorem statement, 
%%we can find the characteristic function $\Phi_X$ evaluated at $u \in \R^n$ as 
%%\EQ{
%%\Phi_X(u) 
%%= \prod_{x\in\R^n}dxe^{j\inner{u,x}}f_X(x)dx
%%= .
%%}
%%
%%To find the density of $X$ in terms of density of $Z$, we use the transformation of random vectors. 
%%Therefore, we have 
%%\EQ{
%%P\Big(\cap_{i=1}^n\set{X_i \le x_i}\Big) 
%%= P\Big(\cap_{i=1}^n\set{\sum_{k=1}^nA_{i,k}Z_k \le x_i-\mu_i}\Big)
%%}
%\end{rem}
%%\begin{shaded}
%%\begin{exmp}
%%Let $X:\Omega\to\R^n$ be an \iid Gaussian random vector with the mean $\mu_1 \triangleq \E X_1$ and variance $\sigma^2 = \E(X_1-\mu_1)^2$ for each of the component random variables. 
%%Then, we have the joint density
%%\EQ{
%%f_X(x) 
%%= \prod_{i=1}^nf_{X_1}(x_i) 
%%= \frac{1}{(2\pi\sigma^2)^{n/2}}\exp(-\sum_{i=1}^n\frac{(x-\mu_1)^2}{2\sigma^2}).
%%}  
%%This implies that $\mu = \mu_1\bU$ and $\Sigma = \sigma^2I$.
%%\end{exmp}
%%\end{shaded}
%\begin{lem} 
%If the components of the Gaussian random vector are uncorrelated, 
%then they are independent. 
%\end{lem} 
%\begin{proof}
%If a Gaussian vector is uncorrelated, then the covariance matrix $\Sigma$ is diagonal. 
%It follows that we can write $f_X(x) = \prod_{i=1}^nf_{X_i}(x_i)$ for all $x\in\R^n$.  
%\end{proof}
%%\begin{lem}
%%A Gaussian random vector $X: \Omega \to \R^n$ with mean $\mu$ and covariance $\Sigma$ is \iid Gaussian iff $\Sigma = \sigma^2I$.
%%\end{lem}
%%\begin{lem} 
%%The characteristic function $\Phi_X: \R^n \to \C$ for a Gaussian random vector $X: \Omega \to \R^n$ with mean $\mu$ and covariance $\Sigma$ is given by 
%%\EQ{
%%\Phi_X(u) = .
%%}
%%\end{lem}
%%\begin{proof}
%%Since $X: \Omega \to \R^n$ is a random Gaussian vector with mean $\mu$ and covariance $\Sigma$, 
%%random variable $\inner{u,X} = \sum_{i=1}^nu_iX_i$ is Gaussian with mean $\E\inner{u,X} = \inner{u,\mu}$ and covariance $u^T\Sigma u$. 
%%%$\sum_{i,j=1}^nu_i\E[(X_i-\mu_i)(X_j-\mu_j)]u_j = u^T\E(X-\mu)(X-\mu)^T}u = u^T\Sigma u$.  
%%\end{proof}


%%%%%%%%%%%%%%%%%%%%%%%%%%%%%%%%%%%%%%%%%%%%%%%%%%%%%%%%%%%%%%%%%%
\section{Gaussian Random Vectors}
%%%%%%%%%%%%%%%%%%%%%%%%%%%%%%%%%%%%%%%%%%%%%%%%%%%%%%%%%%%%%%%%%%

\begin{defn}
For a random vector $X: \Omega\to \R^n$ defined on a probability space $(\Omega,\sF,P)$, 
we can define the characteristic function $\Phi_X:\R^n\to\C$ by $\Phi_X(u) \triangleq \E e^{j \inner{u,X}}$ where $u \in \R^n$. 
\end{defn}
\begin{rem}
If $X:\Omega\to\R^n$ is an independent random vector, 
then $\Phi_X(u) = \prod_{i=1}^n\Phi_{X_i}(u_i)$ for all $u \in \R^n$.
\end{rem}
\begin{defn}%[Gaussian random vectors] 
For a probability space $(\Omega, \sF, P)$, 
\textbf{Gaussian random vector} is a continuous random vector $X: \Omega \to \R^n$ defined by its density function
\EQ{
f_X(x) \triangleq \frac{1}{\sqrt{(2\pi)^n\det(\Sigma)}}\exp\left(-\frac{1}{2}(x-\mu)^T\Sigma^{-1}(x-\mu)\right)\text{ for all }x \in \R^n,
}
where the vector $\mu \in \R^n$ and the positive definite matrix $\Sigma \in \R^{n \times n}$. 
\end{defn}
%\begin{rem} 
%We can verify that for a Gaussian vector, the mean vector is $\E X = \mu$ and the covariance matrix is $\E(X-\mu)(X-\mu)^T = \Sigma$.
%\end{rem}
\begin{rem} 
For a Gaussian random vector with vector $\mu = (\mu_1, \dots, \mu_1)$ for some real scalar $\mu_1$ and matrix $\Sigma = \sigma^2I$ for some positive $\sigma^2 \in \R_+$, 
we can write its density as 
$
f_X(x) = \frac{1}{(2\pi\sigma^2)^{n/2}}\exp\left(-\frac{1}{2}\sum_{i=1}^n\frac{(x_i-\mu_1)^2}{\sigma^2}\right)\text{ for all }x \in \R^n.
$
It follows that $X$ is an \iid random vector with each component being a Gaussian random variable with mean $\mu_1$ and variance $\sigma^2$. 
The characteristic function $\Phi_X$ of an \iid Gaussian random vector $X:\Omega\to \R^n$ parametrized by $(\mu_1,\sigma^2)$ is given by 
$
\Phi_X(u) = \prod_{i=1}^n\Phi_{X_i}(u_i) = \exp\Big(-\frac{\sigma^2}{2}\sum_{i=1}^nu_i^2 +j\mu_1\sum_{i=1}^nu_i\Big).
$
\end{rem}
\begin{lem} 
For an \iid zero mean unit variance Gaussian vector $Z:\Omega\to\R^n$, vector $\alpha \in \R^n$, and scalar $\mu \in \R$, 
the affine combination $Y \triangleq \mu + \inner{\alpha, Z}$ is a Gaussian random variable.  
\end{lem}
\begin{proof}
From the linearity of expectation and the fact that $Z$ is a zero mean vector, 
we get $\E Y = \mu$. 
Further, from the linearity of expectation and the fact that $\E[ZZ^T] = I$, we get 
\EQ{
\sigma^2 \triangleq \Var(Y) = \E(Y-\mu)^2 = \sum_{i=1}^n\sum_{k=1}^n\alpha_i\alpha_k\E[Z_iZ_k] = \inner{\alpha,\alpha} = \norm{\alpha}^2_2 = \sum_{i=1}^n\alpha_i^2.
}
To show that $Y$ is Gaussian, 
it suffices to show that $\Phi_Y(u) = \exp(-\frac{u^2\sigma^2}{2}+ju\mu)$ for any $u \in \R$. 
Recall that $Z$ is an independent random vector with individual components being identically zero mean unit variance Gaussian. 
Therefore, $\Phi_{Z_i}(u) = \exp(-\frac{u^2}{2})$, 
and we can compute the characteristic function of $Y$ as 
\EQ{
\Phi_Y(u) 
= \E e^{juY} 
= e^{ju\mu}\E \prod_{i=1}^ne^{ju\alpha_iZ_i}
= e^{ju\mu}\prod_{i=1}^n\Phi_{Z_i}(u\alpha_i)
= \exp(-\frac{u^2\sigma^2}{2}+ju\mu).
}
\end{proof}
\begin{thm}%[Gaussian random vector] 
%Let $Z:\Omega\to\R^n$ is an \iid zero mean unit variance Gaussian random vector, then $X \triangleq \mu + \Sigma^{\frac{1}{2}}Z$ is a Gaussian random vector with parameters $(\mu, \Sigma)$. 
A random vector $X: \Omega\to\R^n$ is Gaussian with parameters $(\mu, \Sigma)$ iff $Z \triangleq \Sigma^{-\frac{1}{2}}(X-\mu)$ is an \iid zero mean unit variance Gaussian random vector. 
%mean vector $\mu = \E X \in \R^n$, 
%and covariance matrix $\Sigma \triangleq \E(X-\mu)(X-\mu)^T$. 
\end{thm}
\begin{proof}
Let  $X = \mu + \Sigma^{\frac{1}{2}}Z$ for an \iid zero mean unit variance Gaussian random vector $Z:\Omega\to\R^n$, 
then we will show that $X$ is a Gaussian random vector by transformation of random vector densities. 
%Since $Z = \Sigma^{-\frac{1}{2}}(X-\mu)$, 
Since the $(i,j)$th component of the Jacobian matrix $J(x)$ is given by 
$J_{ij}(x) = \frac{\partial x_j}{\partial z_i} = \Sigma^{\frac{1}{2}}_{i,j}$ for all $i,j\in [n],$
we can write the Jacobian matrix $J(x) = \Sigma^{\frac{1}{2}}$, 
Since the density of $Z$ is $f_Z(z) = \frac{1}{\sqrt{(2\pi)^n}}\exp(-\frac{1}{2}z^Tz)$, 
and from the transformation of random vectors, we get 
\EQ{
f_X(x) 
= \frac{f_Z(\Sigma^{-\frac{1}{2}}(x-\mu))}{\det(\Sigma^{\frac{1}{2}})}
= \frac{1}{(2\pi)^{n/2}\det(\Sigma)^{1/2}}\exp\Big(-\frac{1}{2}(x-\mu)^T\Sigma^{-1}(x-\mu)\Big),\quad x\in \R^n.
}  

Conversely, we can show that if $X$ is a Gaussian random vector, then $Z = \Sigma^{-\frac{1}{2}}(X-\mu)$ is an \iid zero mean unit variance Gaussian random vector by transformation of random vectors. 
%Since $X = \mu + \Sigma^{\frac{1}{2}}Z$, it follows that the mean vector of $X$ is $\mu = \E X$ and the covariance matrix is $\Sigma = \E(X-\mu)(X-\mu)^T.$
%It is easy to verify that the mean vector $\E X = \mu$ and the covariance matrix is $\E(X-\mu)(X-\mu)^T = \Sigma^{\frac{1}{2}}(\E ZZ^T)\Sigma^{\frac{1}{2}} = \Sigma.$
\end{proof}
\begin{rem}
A random vector $X:\Omega\to\R^n$ with mean $\mu \in \R^n$ and covariance matrix $\Sigma \in \R^{n\times n}$ is Gaussian iff $X$ can be written as 
$
X = \mu + \Sigma^{\frac{1}{2}}Z,
$
for an  \iid Gaussian random vector $Z: \Omega \to \R^n$ with mean~$0$ and variance~$1$. 
It follows that $\E X = \mu$ and $\Sigma = \E(X-\mu)(X-\mu)^T$. 
\end{rem}
\begin{rem}
We observe that the components of the Gaussian random vector $X = \mu + A Z$ for $A = \Sigma^{\frac{1}{2}}$ are Gaussian random variables with mean $\mu_i$ and variance $\sum_{k=1}^nA_{i,k}^2 = (AA^T)_{i,i} = \Sigma_{i,i}$, 
since each component $X_i = \mu_i + \sum_{k=1}^nA_{i,k}Z_k$ is an affine combination of zero mean unit variance \iid random variables. 
\end{rem}
\begin{rem}
For any $u \in \R^n$, 
we compute the characteristic function $\Phi_X$ from the distribution of $Z$ as 
\EQ{
\Phi_X(u) 
= \E e^{j\inner{u,X}} 
= \E\exp\Big(j\inner{u,\mu} + j\inner{A^Tu,Z}\Big)
= \exp(j\inner{u,\mu})\Phi_Z(A^Tu) 
= \exp(j\inner{u,\mu}-\frac{1}{2}u^T\Sigma u).
}
\end{rem}
\begin{lem} 
If the components of the Gaussian random vector are uncorrelated, 
then they are independent. 
\end{lem} 
\begin{proof}
If a Gaussian vector is uncorrelated, then the covariance matrix $\Sigma$ is diagonal. 
It follows that we can write $f_X(x) = \prod_{i=1}^nf_{X_i}(x_i)$ for all $x\in\R^n$.  
\end{proof}
\end{document}